\section{Introduction}
\label{sec:intro}

Generative retrieval replaces nearest-neighbor search over dense embeddings with autoregressive generation of a discrete identifier that points directly to the target candidate~\cite{tay2022dsi} (see~\cite{li2024survey} for a survey). GENIUS~\cite{kim2025genius} applies this to universal multimodal search: a residual quantizer (RQ)~\cite{lee2022rqvae} assigns every candidate a short sequence of discrete codes, and a decoder learns to generate a query's target code sequence directly, retrieving in time independent of corpus size. Since the code sequence is the only channel through which a candidate's identity survives into the decoder, codebook quality is a ceiling on what the rest of the system can achieve.

A natural lever for improving it is \emph{balance regularization}: an auxiliary loss discouraging a small subset of codes from absorbing most assignments. Whether this helps is not obvious a priori: broader usage should reduce collisions, but it pulls the encoder from the pure-reconstruction optimum, and its effect may depend on task structure -- a diverse pool may benefit from spreading code usage, while a pool of near-duplicates may instead need the encoder to preserve the subtle differences any diversity pressure would blur.

We first reproduce GENIUS end-to-end from its officially released Stage-1 checkpoint (Section~\ref{sec:setup}), before studying any regularization variant. We then test balance regularization across three regimes -- MSCOCO and VisualNews (broad-domain captioning, two different image/text distributions) and FashionIQ (fine-grained composed fashion retrieval) -- training RQ tokenizers under multiple balance strengths per dataset, with a multi-seed check on every headline number that has already caught one single-run artifact in this exact study (Section~\ref{sec:limitations}), including the decoupling experiment of Section~\ref{sec:direction-aware}, whose constrained-decoding evaluation initially failed on every retrain attempt for reasons we root-cause, fix, and then verify at 3 seeds (Section~\ref{sec:limitations}). Concretely, we ask:

\begin{itemize}
\item \textbf{RQ1:} Which ID-structure properties (codebook utilization, collision rate, reconstruction quality) track downstream retrieval, and is that relationship consistent across retrieval regimes? (Section~\ref{sec:rq1})
\item \textbf{RQ2:} Does balance regularization improve downstream Recall, and at what strength? (Section~\ref{sec:rq2rq3})
\item \textbf{RQ3:} Is the answer stable -- across retrieval direction, across datasets, and under a per-modality decoupling of the regularizer? (Sections~\ref{sec:rq2rq3}--\ref{sec:generality})
\end{itemize}

\paragraph{Contribution.} On MSCOCO, the regularization that \emph{significantly improves} text-to-image retrieval simultaneously destroys image-to-text retrieval on the same dataset. This holds on two independent axes of verification: 5 decoder seeds on the original tokenizer (non-overlapping distributions, $p<10^{-6}$ -- a conclusion that itself survived a scare, Section~\ref{sec:limitations}), and 3 fully independent RQ tokenizers retrained from scratch (every regularized tokenizer outperforms every unregularized one on text-to-image, and the image-to-text collapse holds on all three). We localize this collapse directly to the tokenizer using a dense, decoder-free probe, and partially explain it via text embeddings' greater anisotropy under regularization. Testing generality on VisualNews shows the collapse reproduces cleanly on a second broad-domain corpus, while the text-to-image effect \emph{reverses} -- strong regularization there substantially \emph{hurts} T$\to$I instead of helping it (a $53\%$ relative drop, seed-stable but not significance-tested, unlike MSCOCO's result) -- showing the collapse is an architectural property while the T$\to$I effect is dataset-specific in a way the broad/fine-grained axis alone does not capture. FashionIQ shows any regularization hurts fine-grained retrieval outright. So this study's one confirmed retrieval benefit is real but narrow: significant and tokenizer-verified on MSCOCO, reliably absent or reversed everywhere else we tested.

We then test the natural fix to the trade-off -- decoupling regularization strength per target modality -- at three settings, including regularizing only the image side and, symmetrically, only the text side. It fails on all three: image-to-text retrieval stays collapsed even at zero text-side loss weight, collapses just as completely when text-side weight alone is nonzero, and text-to-image retrieval becomes \emph{worse} than either baseline in every asymmetric setting. This is architectural, not a tuning failure -- image and text rows share one encoder and one EMA-updated codebook, so a modality's loss weight cannot shield its code structure from the other's regularization -- supported by direct evidence of cross-modal code restructuring under zero cross-modal loss weight, and corroborated at the tokenizer level by the same decoder-free probe used to localize the main collapse.

This paper reproduces GENIUS end-to-end, then generalizes and analyzes one of its core design choices across regimes it was never evaluated under. To our knowledge, no prior work has localized an image-to-text collapse in a generative multimodal retriever to the tokenizer rather than the decoder with a direct probe, replicated across two broad-domain corpora, where the same probe also shows VisualNews's text-to-image \emph{drop} is tokenizer-level too -- neither direction is decoder-introduced. Nor has prior work shown the effective regularization strategy for RQ-based semantic IDs is retrieval-regime-dependent in a way that is not simply ``broad vs.\ fine-grained'' -- a single dataset's positive result does not license a general recommendation. And, at 3-seed and tokenizer-level verification (Section~\ref{sec:limitations}), no prior work has shown that per-modality decoupling of a generative retriever's codebook regularization is undermined by the architecture itself, not by tuning.
